\section{Limitations}
Despite the achieved progress, the current project still has important limitations that define the next development stage:

\begin{itemize}[leftmargin=1.5em]
\item \textbf{Dataset and annotation constraints:} model behavior is strongly influenced by the quality, coverage, and annotation consistency of available ISL data.
\item \textbf{Signer and domain variability:} recognition stability can degrade for unseen signer styles, camera viewpoints, rapid motion, or heavy occlusion/background clutter.
\item \textbf{Hardware sensitivity:} low-latency operation depends on compute resources; lower-end devices may face higher delay and reduced throughput.
\item \textbf{Language post-processing limits:} grammar correction currently improves readability but does not yet provide full semantic translation quality in all sentence contexts.
\item \textbf{Long-context understanding limits:} very long or complex continuous signing segments can still challenge temporal boundary handling and decoding confidence.
\item \textbf{Real-world deployment variability:} practical conditions such as illumination changes, camera quality differences, and network constraints can impact live performance.
\item \textbf{Evaluation depth:} current evaluation emphasizes technical and interface validation; broader user-centered field validation across diverse signer groups remains to be expanded.
\end{itemize}

These limitations do not invalidate the system; instead, they indicate where targeted engineering and data improvements are required to move from a strong prototype baseline to robust large-scale assistive adoption.
